% ---------------------------------------------------------------------------
%  10-limits.tex -- boundaries, negative results, and threats to validity
% ---------------------------------------------------------------------------
\section{Limits, negative results, and threats to validity}
\label{sec:limits}

\subsection{Where the design does not scale, by construction}

\begin{itemize}[leftmargin=1.2em, itemsep=3pt]
\item \textbf{One discrete state, many zones.}  A model whose reachable states
      all share one discrete state has one stripe, hence one lock and one
      antichain, and does not parallelise (\code{single}, \cref{fig:scaling}).
      This is inherent to sharding by discrete key.  A design that also sharded
      \emph{within} an antichain would need a different concurrency argument
      than \cref{thm:coverage}, because the merge in step (4) is global to the
      chain.

\item \textbf{Narrow levels.}  Level-synchronous BFS cannot use more workers
      than the current level has states.  Models with long, thin search graphs
      pay the barrier on every level and gain nothing.

\item \textbf{Diagonal constraints.}  These disable $\aLU$ entirely
      \cite{Bouyer04}; the engine falls back to exact inclusion with a
      diagonal-preserving widening, and imposes a state cap because zones can
      otherwise grow without bound.  The parallel machinery is unaffected, but
      the algorithmic advantage over \uppaal{} disappears.

\item \textbf{Liveness is still sequential.}  \code{E[] phi} and \code{p --> q}
      build a zone graph and then run a cycle detection over it; only the
      \code{E<>} family (including \code{E<> deadlock}) runs on the parallel
      engine.  This is a real gap and should be stated as such rather than
      glossed over.

\item \textbf{The discrete key is $63$ bits.}  Locations and integer variables
      are bit-packed into one word; wide networks are rejected with an explicit
      breakdown of which fields overflowed.  This is a deliberate trade --- the
      packed key is what makes hashing and comparison single-instruction --- but
      it is a limit on model size in a dimension unrelated to state count.
\end{itemize}

\subsection{Things that were tried and removed}

These are worth a paragraph in the paper.  They are evidence that the design
space was actually explored, and each one closes off an obvious reviewer
suggestion.

\begin{itemize}[leftmargin=1.2em, itemsep=3pt]
\item \textbf{Per-thread local antichains with a drain at the barrier.}  Each
      worker inserts into a private table during the level and an owner drains
      them into the shared one at the boundary, so expansion takes no lock at
      all.  Correct, and consistently \emph{slower}: at $8$ threads,
      \code{phil11} $1.48\times$, \code{viking12} $1.29\times$,
      \code{twolock10} $1.90\times$, \code{fischer10} $1.92\times$ worse.  The
      reason is precise and instructive: removing the lock saves the $\approx
      15\%$ of thread time that was spin, but a private antichain cannot see the
      shared coverage probe that answers $83\%$ of successors, so the drain
      pays back more than the lock cost.  On \code{phil12} the drain handled
      $3.2$\,M gathered pairs, $1.5$\,M of them exact duplicates, spending
      $19.4$\,s of thread time against a $3.8$\,s wall.  \emph{A local-buffer
      design is only viable with a shared coverage probe.}

\item \textbf{Expansion ordering by generality.}  A zone that a later, more
      general zone subsumes was expanded for nothing ($45\%$ of arena nodes on
      \code{phil12} end up retired), so expanding the most general zones of a
      level first ought to let them kill their narrower siblings a level
      earlier.  The generality proxy was free (a minimal constraint system
      stores a more general zone in a \emph{shorter} record) and ordering a
      worklist cannot affect a fixpoint, so it was sound.  It does not pay:
      \code{phil12} $0.95\times$ and \emph{more} states ($841\,637$ vs
      $803\,845$), \code{fischer11} $0.98\times$, \code{twolock11}
      $1.00\times$, \code{viking15} $0.98\times$.  The retirement is driven by
      zones arriving levels apart, which no intra-level ordering can reach.

\item \textbf{Test-and-test-and-set with exponential back-off}, and
      \textbf{writer preference}, in the stripe lock.  Neither helped; back-off
      was slightly worse (a waiter that probes less often misses the release
      window) and writer preference merely moved wait from writers to readers,
      because the stripes are genuinely serialised rather than unfairly
      scheduled.
\end{itemize}

\subsection{Threats to validity}

\begin{itemize}[leftmargin=1.2em, itemsep=3pt]
\item \textbf{The headline ratio compares a $16$-thread tool to a $1$-thread
      tool.}  It is a fair engineering comparison (it is what a user gets) but
      it is not an algorithmic one.  The state-count columns are what separate
      the two effects, and they should be given equal prominence.

\item \textbf{Missing sequential baseline.}  Without $1$--$8$ thread data the
      absolute parallel speedup is not established, only the $16 \to 32$ slope.
      This is the most important gap to close before submission.

\item \textbf{Benchmark selection.}  Five families, all classical, all
      process-replicated.  Replicated networks are exactly the shape this design
      is best at (many discrete states, short antichains).  \code{single} is
      included as the adversarial case, but one adversarial family is thin
      evidence; a second, e.g.\ a model with wide integer state and few
      locations, would strengthen the claim.

\item \textbf{Merge order affects reported state counts.}  Two runs at $32$
      threads can report different counts (\cref{thm:determinism}).  Any table
      of state counts must therefore either be produced at one thread or state
      the variation.

\item \textbf{The oracle is \code{verifyta} itself.}  Differential testing
      against the reference tool establishes agreement, not absolute
      correctness.  The compiled-in invariant checkers (\cref{sec:eval}) are
      the independent evidence, and they check the specific lemmas of
      \cref{sec:theory} rather than the end-to-end verdict.
\end{itemize}
